\chapter{Grafy i węzły w konstrukcjach czterowymiarowych}
\label{ch:grafy-wezly-cztery}

Rozdział~\ref{ch:freedman-formy} oblicza formę przecięcia
z obramowanego splotu, a Rozdziały~\ref{ch:freedman-klasyfikacja}
i~\ref{ch:egzotyczne-r4} rozróżniają kategorie TOP i DIFF.
Teraz zbudujemy dwie konkretne operacje na \emph{zwartych}
czterorozmaitościach. Pierwsza zastępuje konfigurację sfer
wymierną kulą homologii; druga wymienia otoczenie torusa
na iloczyn okręgu z zewnętrzem węzła. Pokażemy rachunki
wynikające z grafu i ciągów dokładnych. Istnienie
specjalnego wypełnienia oraz formułę Seiberga--Wittena
przyjmiemy jako wyraźnie wskazane wyniki zewnętrzne.

\section{Od ważonego drzewa do formy przecięcia}
\label{sec:plumbing-graf-cztery}
\index{plumbing!grafowy}

Niech $\Gamma$ będzie skończonym drzewem o wierzchołkach
$v_1,\ldots,v_r$ z wagami $a_i\in\Z$. Wierzchołkowi
$v_i$ przypisujemy zorientowaną wiązkę dysków
$D^2\to E(a_i)\to S^2$ o liczbie Eulera $a_i$.
Dla każdej krawędzi wybieramy w obu wiązkach osobne
bloki $D^2_{\text{baza}}\times D^2_{\text{włókno}}$ i sklejamy
je przez zamianę czynników. Wybieramy orientacje tak,
by sekcje zerowe przecinały się w każdym sklejanym
bloku z lokalną liczbą $+1$. Po wygładzeniu narożników
otrzymujemy gładką czterorozmaitość z brzegiem
$P_\Gamma$, zwaną \emph{plumbingiem grafowym}.

\begin{stwierdzenie}[Macierz plumbingowa i homologia brzegu]
\label{prop:plumbing-macierz-cztery}
Rozmaitość $P_\Gamma$ jest po prostu spójna,
$H_2(P_\Gamma;\Z)\cong\Z^r$, a klasy sekcji zerowych
$S_i$ tworzą bazę. Macierz formy przecięcia ma wpisy
\[
 Q_\Gamma(i,j)=
 \begin{cases}
 a_i,&i=j,\\
 1,&v_i,v_j\text{ są połączone krawędzią},\\
 0,&\text{w pozostałych przypadkach}.
 \end{cases}
\]
Ponadto $H_1(\partial P_\Gamma;\Z)
\cong\operatorname{coker}Q_\Gamma$.
\end{stwierdzenie}
\begin{proof}
Każda wiązka $E(a_i)$ retrahuje się na sekcję zerową
$S^2$. Kolejne bloki plumbingowe są ściągalne;
ponieważ graf sklejeń jest drzewem, van Kampen daje
$\pi_1(P_\Gamma)=1$, a Mayer--Vietoris dodaje po
jednym wolnym generatorze $H_2$ przy każdym nowym
wierzchołku. Sekcja zerowa ma samoprzecięcie równe
liczbie Eulera jej wiązki normalnej, czyli $a_i$.
Dwie sekcje spotykają się w jednym dodatnim punkcie
dokładnie wtedy, gdy ich wierzchołki są połączone.
To daje macierz $Q_\Gamma$.

Niech $Y=\partial P_\Gamma$. Dualność
Poincarégo--Lefschetza i twierdzenie o współczynnikach
uniwersalnych, przy $H_1(P_\Gamma)=0$, identyfikują
$H_2(P_\Gamma,Y;\Z)$ z
$\operatorname{Hom}(H_2(P_\Gamma;\Z),\Z)$.
W długim ciągu pary odcinek
\[
 H_2(P_\Gamma;\Z)\xrightarrow{\ j\ }
 H_2(P_\Gamma,Y;\Z)\longrightarrow H_1(Y;\Z)
 \longrightarrow H_1(P_\Gamma;\Z)=0
\]
ma mapę $j$ reprezentowaną przez parowanie przecięcia,
a więc przez $Q_\Gamma$. Jego kokernel jest $H_1(Y;\Z)$.
\end{proof}

Sam wyznacznik $Q_\Gamma$ daje tylko rząd grupy
$H_1(Y)$, gdy macierz jest nieosobliwa. Nie rozpoznaje
typu dyfeomorficznego brzegu ani jego orientacji.
Graf $E_8$ z Rysunku~\ref{fig:freedman-E8-graph}
jest szczególnym przykładem tej konstrukcji,
po zmianie orientacji części sfer odwracającej
znaki wpisów przy krawędziach drzewa.

\section{Łańcuch \texorpdfstring{$C_p$}{Cp} i wymierna kontrakcja}
\label{sec:kontrakcja-wymierna-cztery}
\index{kontrakcja wymierna}

Niech $p\geq2$. Definiujemy $C_p=P_{\Gamma_p}$ dla
łańcucha $p-1$ sfer z wagami
$-(p+2),-2,\ldots,-2$; dwójka występuje $p-2$ razy.
Jego macierz $Q_p$ ma jedynki bezpośrednio obok
przekątnej. Dla $p=2$ jest to pojedyncza sfera
o samoprzecięciu $-4$.

\begin{figure}[htbp]
\centering
\begin{tikzpicture}[font=\small]
 \node[circle,draw=manifoldred!80!black,fill=manifoldred!12,
       minimum size=11mm] (a) at (0,0) {$-(p+2)$};
 \node[circle,draw=manifoldblue!80!black,fill=manifoldblue!12,
       minimum size=9mm] (b) at (2.1,0) {$-2$};
 \node[circle,draw=manifoldblue!80!black,fill=manifoldblue!12,
       minimum size=9mm] (c) at (3.7,0) {$-2$};
 \node (dots) at (5.1,0) {$\cdots$};
 \node[circle,draw=manifoldblue!80!black,fill=manifoldblue!12,
       minimum size=9mm] (d) at (6.4,0) {$-2$};
 \draw[manifoldblue!80!black,thick] (a)--(b)--(c)--(dots)--(d);
 \node[below=5mm of c] {schemat dla $p\geq6$; razem $p-1$ wierzchołków};
\end{tikzpicture}
\caption{Ważony graf konfiguracji $C_p$ Fintushela--Sterna.}
\label{fig:graf-cp-cztery}
\end{figure}

\begin{stwierdzenie}[Forma i brzeg konfiguracji]
\label{prop:cp-rachunek-cztery}
Macierz $Q_p$ jest ujemnie określona,
$\det Q_p=(-1)^{p-1}p^2$ oraz
$H_1(\partial C_p;\Z)\cong\Z/p^2\Z$.
Przy konwencji, w której $L(n,q)$ oznacza
$n/q$-chirurgię na niewęźle, zachodzi
$\partial C_p\cong L(p^2,1-p)$ z orientacją brzegową.
\end{stwierdzenie}
\begin{proof}
Niech $d_k$ będzie wyznacznikiem początkowego bloku
$k\times k$ macierzy $-Q_p$. Mamy $d_0=1$,
$d_1=p+2$ oraz, po rozwinięciu trójdiagonalnego
wyznacznika, $d_k=2d_{k-1}-d_{k-2}$.
Indukcja daje $d_k=1+k(p+1)$; w szczególności
$d_{p-1}=1+(p-1)(p+1)=p^2$.
Wszystkie początkowe minory główne $-Q_p$ są dodatnie,
więc kryterium Sylvestera daje ujemną określoność $Q_p$.

Poprzednie stwierdzenie identyfikuje homologię brzegu
z $\operatorname{coker}Q_p$. Sama wartość wyznacznika
daje rząd $p^2$. Aby ustalić strukturę grupy, oznaczmy
generatory kokernela przez $e_1,\ldots,e_{p-1}$.
Pierwsza relacja daje $e_2=(p+2)e_1$; następne
$e_{j+1}=2e_j-e_{j-1}$. Indukcyjnie
$e_j=[1+(j-1)(p+1)]e_1$.
Ostatnia relacja staje się $p^2e_1=0$.
Wszystkie generatory są więc wielokrotnościami $e_1$,
a rząd grupy wynosi $p^2$; dla $p=2$ wynika to
bezpośrednio z $Q_2=[-4]$.

Współczynnik wymierny liniowego łańcucha chirurgii
jest ujemnym ułamkiem łańcuchowym
$-[p+2,2,\ldots,2]^-$.
Ponieważ łańcuch $m$ dwójek ma wartość $(m+1)/m$,
otrzymujemy dla $p\geq3$
\[
 -[p+2,\underbrace{2,\ldots,2}_{p-2}]^-
 =-\left(p+2-\frac{p-2}{p-1}\right)
 =-\frac{p^2}{p-1}=\frac{p^2}{1-p}.
\]
Dla $p=2$ ten sam wzór daje $-4$.
Standardowy rachunek chirurgii liniowego splotu Hopfa
identyfikuje jego brzeg z przestrzenią soczewkową
o tym współczynniku i ustalonej orientacji.
Ten ostatni krok wykorzystuje regułę eliminacji
składowych ramowanego splotu, a nie sam wyznacznik.
\end{proof}

\begin{twierdzenie}[Casson--Harer; wynik zewnętrzny]
\label{thm:kula-bp-cztery}
Dla każdego $p\geq2$ istnieje gładka zwarta
czterorozmaitość $B_p$, której orientację można
dobrać tak, by
$\partial B_p\cong L(p^2,1-p)$,
$H_*(B_p;\Q)\cong H_*(B^4;\Q)$ oraz
$\pi_1(B_p)\cong\Z/p\Z$.
W modelu uchwytowym homomorfizm
$\pi_1(\partial B_p)\to\pi_1(B_p)$ jest surjektywny.
\end{twierdzenie}
\noindent Konstrukcja $B_p$ i identyfikacja jej brzegu
są nietrywialnym wejściem zewnętrznym; rachunek formy
$Q_p$ sam ich nie dowodzi.

Jeśli $C_p$ jest gładko zanurzona w zamkniętej
zorientowanej czterorozmaitości $X$, kładziemy
$W=X\setminus\operatorname{int}C_p$ i, po ustaleniu
zgodnej z orientacjami identyfikacji brzegów,
\[
 X_{(p)}=W\cup_{\partial C_p}B_p.
\]
Tę operację nazywamy \emph{wymierną kontrakcją}.
Jej nazwa odnosi się do homologii z $\Q$;
$B_p$ nie jest zwykłą kulą, bo $H_1(B_p;\Z)=\Z/p\Z$.

\begin{stwierdzenie}[Zmiana niezmienników]
\label{prop:kontrakcja-niezmienniki-cztery}
Przy powyższych założeniach
\[
 \chi(X_{(p)})=\chi(X)-(p-1),\qquad
 \sigma(X_{(p)})=\sigma(X)+(p-1).
\]
Ponadto $b_2^+$ pozostaje bez zmiany, a $b_2^-$
maleje o $p-1$. Jeśli $\pi_1(W)=1$, to również
$\pi_1(X_{(p)})=1$.
\end{stwierdzenie}
\begin{proof}
Brzeg $Y=\partial C_p$ jest wymierną sferą homologiczną,
więc ciąg Mayera--Vietorisa nad $\Q$ daje
$H_2(X;\Q)\cong H_2(W;\Q)\oplus H_2(C_p;\Q)$
oraz $H_2(X_{(p)};\Q)\cong H_2(W;\Q)$.
Pierwsza suma jest ortogonalna dla przecięć klas
reprezentowanych we wnętrzach części. Usuwamy więc
ujemnie określoną formę rangi $p-1$.
Stąd wzory dla sygnatury i $b_2^\pm$.
Plumbing ma homologię $p-1$ sfer i punktu,
zatem $\chi(C_p)=p$; dla wymiernej kuli
$\chi(B_p)=1$. Addytywność charakterystyki Eulera
przy wspólnym brzegu daje pierwszy wzór.
Wreszcie van Kampen dla $W\cup_Y B_p$ przy
$\pi_1(W)=1$ dzieli $\pi_1(B_p)$ przez normalny
obraz $\pi_1(Y)$. Ten obraz jest całą grupą na mocy
Twierdzenia~\ref{thm:kula-bp-cztery}.
\end{proof}

Prosta spójność $X$ sama nie zastępuje warunku
$\pi_1(W)=1$. W zastosowaniach trzeba obliczyć
grupę dopełnienia konfiguracji. Obecny rachunek
nie rozpoznaje też typu gładkiego $X_{(p)}$;
do tego potrzeba dodatkowych niezmienników.

\section{Spinory i niezmiennik Seiberga--Wittena}
\label{sec:sw-wprowadzenie-cztery}
\index{Seiberg--Witten!niezmiennik}
\index{struktura Spin-c}

Pola spinorowe i operator Diraca poznaliśmy w
Rozdziałach~\ref{ch:grass-spin}
i~\ref{ch:hodge-bochner}. Grupa
$\mathrm{Spin}^c(4)$ jest ilorazem
$(\mathrm{Spin}(4)\times U(1))/\{(1,1),(-1,-1)\}$.
Struktura $\mathrm{Spin}^c$ to podniesienie wiązki
zorientowanych ramek ortonormalnych przez jej naturalne
odwzorowanie na $SO(4)$. Na zorientowanej gładkiej
czterorozmaitości taka struktura daje
wiązki spinorów $S^+,S^-$ oraz zespoloną liniową
wiązkę wyznacznikową $L$, odpowiadającą charakterowi
$[s,z]\mapsto z^2$. Jej klasa
$c_1(L)\in H^2(X;\Z)$ redukuje się modulo $2$
do $w_2(X)$. Dla połączenia $A$ w $L$ i spinora
$\psi\in\Gamma(S^+)$ równania Seiberga--Wittena
mają, przy jednej z konwencji normalizacji, postać
\[
 D_A\psi=0,\qquad F_A^+=q(\psi).
\]
Tu $F_A^+$ jest samodualną częścią krzywizny,
a $q$ kwadratowym odwzorowaniem spinorów.
Po uwzględnieniu grupy cechowania i perturbacji
ogólnej pozycji, dla $b_2^+(X)>1$ oraz zerowego
wymiaru oczekiwanego przestrzeni rozwiązań,
jej zorientowane zliczenie definiuje liczbę
$SW_X(\mathfrak s)$. Wymiar oczekiwany wynosi
\[
 d(\mathfrak s)=
 \frac{c_1(L)^2-(2\chi(X)+3\sigma(X))}{4}.
\]
Ten wzór wymaga twierdzenia indeksowego, a zwartość,
regularność i niezmienniczość zliczenia wymagają
analizy równań; nie wynikają z poprzednich rachunków
homologii. Będziemy używać tylko ustalonych w literaturze
wartości i formuły chirurgicznej. Przy $b_2^+=1$
pojawiają się komory, więc bezkomorowej postaci
formuły poniżej tam nie stosujemy.

\section{Chirurgia wzdłuż węzła i egzotyczne K3}
\label{sec:chirurgia-wezel-cztery}
\index{chirurgia wzdłuż węzła}
\index{K3!egzotyczne struktury}

Niech $T\subset X$ będzie gładkim torusem o
samoprzecięciu zero i wybranym produktowym
otoczeniu $\nu T\cong T^2\times D^2$. Dla węzła
$K\subset S^3$ weźmy jego zewnętrze $E_K$
z Sekcji~\ref{sec:wezly-diagramy-wirtinger}.
Na brzegu $S^1\times E_K$ mamy trzy kierunki:
czynnik $S^1$, meridian $\mu_K$ oraz długość
preferowaną $\lambda_K$. \emph{Chirurgia wzdłuż węzła}
Fintushela--Sterna zastępuje $T^2\times D^2$ przez
$S^1\times E_K$, sklejając $\lambda_K$ z meridianem
usuwanego torusa, a dwa pozostałe kierunki z bazą $T$:
\[
 X_K=(X\setminus\operatorname{int}\nu T)
       \cup_{T^3}(S^1\times E_K).
\]
Sklejenie z ramowaniem i orientacją jest częścią
danych operacji. Dla węzła trywialnego dołączana
część jest ponownie $T^2\times D^2$.

Sprawdźmy zakres elementarnych wniosków.
Zewnętrze węzła ma $H_1(E_K;\Z)=\Z[\mu_K]$,
$H_i(E_K;\Z)=0$ dla $i\geq2$ i
$[\lambda_K]=0$ w $H_1(E_K;\Z)$.
Istotnie, w ciągu Mayera--Vietorisa dla
$S^3=E_K\cup N(K)$, z przecięciem $T^2$,
mapa $H_3(S^3;\Z)\to H_2(T^2;\Z)$ przenosi
klasę fundamentalną na klasę torusa i jest
izomorfizmem. Stąd $H_2(E_K;\Z)=0$;
poprzedni odcinek pokazuje też $H_3(E_K;\Z)=0$.
Następny odcinek ciągu daje izomorfizm
$H_1(T^2;\Z)\cong H_1(E_K;\Z)\oplus H_1(N(K);\Z)$.
Po stronie otoczenia meridian znika, a długość
reprezentuje generator; po stronie zewnętrza
długość jest brzegiem części powierzchni Seiferta.
Zatem meridian daje generator pozostałego składnika.
Wzór Künnetha pokazuje, że $S^1\times E_K$
ma te same grupy homologii co $T^2\times D^2$.
Przy wskazanym sklejeniu mapy z homologii wspólnego
brzegu są zgodne: w stopniu jeden zabijają odpowiednio
meridian torusa i $\lambda_K$, a w stopniu dwa
pozostawiają klasę produktu dwóch pozostałych okręgów.
Ciąg Mayera--Vietorisa w każdym stopniu ma postać
krótkiego ciągu z kokernela mapy brzegowej do
$H_i(X_K;\Z)$ i dalej do jądra mapy w stopniu $i-1$.
Jądro to jest podgrupą wolnej grupy
$H_{i-1}(T^3;\Z)$, zatem jest wolne i ciąg się
rozszczepia. Zgodność map brzegowych daje więc
izomorficzne grupy homologii $X$ i $X_K$.
Obie wymieniane części mają $\chi=0$; ich jedyne
wymierne klasy w $H_2$ są reprezentowane przez
torus o samoprzecięciu zero, więc ich sygnatury
też znikają. Addytywność charakterystyki Eulera
i twierdzenie Nowikowa o addytywności sygnatury
przy sklejeniu wzdłuż $T^3$ dają $\chi(X_K)=\chi(X)$ i
$\sigma(X_K)=\sigma(X)$.
Jeżeli $X\setminus\operatorname{int}\nu T$ jest
po prostu spójne, van Kampen daje $\pi_1(X_K)=1$:
brzeg zabija czynnik $S^1$ i meridian węzła,
a meridian normalnie generuje $\pi_1(E_K)$
na mocy Twierdzenia~\ref{thm:wirtinger-20}.

W następnym twierdzeniu \emph{otoczenie typu cusp}
oznacza standardowe otoczenie kuspidalnego włókna
fibracji eliptycznej. Jako czterorozmaitość z brzegiem
można je opisać śladem dołączenia do $B^4$
uchwytu indeksu dwa wzdłuż trójlistnika z ramowaniem
zero. Torus $T$ ma być regularnym włóknem w tym
otoczeniu. Jest to silniejsze założenie niż samo
$T\cdot T=0$; właśnie pod nim stosujemy formułę SW.

\begin{twierdzenie}[Fintushel--Stern; wynik zewnętrzny]
\label{thm:fs-chirurgia-cztery}
Załóżmy dodatkowo, że $X$ jest zamknięta, zorientowana,
gładka, po prostu spójna, $b_2^+(X)>1$, a homologicznie
nietrywialny torus $T$ leży w otoczeniu typu \emph{cusp}
i ma po prostu spójne dopełnienie. Dla standardowego
sklejenia $X_K$ jest homeomorficzna z $X$.
Jeśli $\mathcal{SW}_X$ oznacza wielomianowy niezmiennik
Seiberga--Wittena w pierścieniu grupowym, to
\[
 \mathcal{SW}_{X_K}=\mathcal{SW}_X\,\Delta_K(u),
 \qquad u=\exp(2[T]),
\]
gdzie $\Delta_K(z)$ jest symetrycznie znormalizowanym
wielomianem Alexandera z Definicji~\ref{def:alexander-wielomian-20},
a $\exp(\beta)$ jest
formalnym symbolem klasy $\beta$ w pierścieniu
grupowym: $\exp(\alpha)\exp(\beta)=\exp(\alpha+\beta)$.
Klasy $[T]$ po obu stronach identyfikujemy przez
standardową korespondencję konstrukcji.
\end{twierdzenie}

Formuła jest głębokim twierdzeniem teorii cechowania.
W szczególności nie wynika z samego zachowania homologii.
Nie wolno jej bez zmiany stosować do kontrakcji wymiernej
przy $b_2^+=1$ ani wnioskować z równości
$\Delta_{K_1}=\Delta_{K_2}$ o dyfeomorfizmie wyników.

\begin{przyklad}[Nieskończenie wiele typów gładkich K3]
\label{prz:ks-k3-cztery}
Powierzchnia K3, oznaczana $E(2)$ jako powierzchnia
eliptyczna, ma $b_2^+=3$ i $\mathcal{SW}_{E(2)}=1$.
Jej regularne włókno $T$ można wybrać w otoczeniu
\emph{cusp} z po prostu spójnym dopełnieniem.
Dla $K_n=T(2,2n+1)$, $n\geq1$, standardowy wzór
na wielomian Alexandera węzła torusowego daje
\[
 \Delta_{K_n}(z)=z^{-n}\frac{z^{2n+1}+1}{z+1}
 =\sum_{j=-n}^{n}(-1)^{n-j}z^j.
\]
Druga równość wynika z dzielenia
$(z^{2n+1}+1)/(z+1)=z^{2n}-z^{2n-1}+\cdots-z+1$.
Stąd $\mathcal{SW}_{E(2)_{K_n}}=\Delta_{K_n}(u)$
ma dokładnie $2n+1$ niezerowych współczynników
przy różnych klasach $2j[T]$. Klasa włókna jest
pierwotna, bo przekrój fibracji przecina je raz.
Liczba klas podstawowych odróżnia więc wyniki
dla różnych $n$, chociaż wszystkie są homeomorficzne
z K3. W przypadku $n=1$ dostajemy trójlistnik
i wielomian $u^{-1}-1+u$.
\end{przyklad}

\par\smallskip
{\footnotesize\noindent\emph{Głębokie źródła:}
\href{https://arxiv.org/abs/hep-th/9411102}
{E.\ Witten, \emph{Monopoles and Four-Manifolds}} ---
równania i indeks wymiaru oczekiwanego;
\href{https://msp.org/pjm/1981/96-1/pjm-v96-n1-p03-s.pdf}
{A.\ J.\ Casson i J.\ L.\ Harer, \emph{Some homology lens
spaces which bound rational homology balls}} --- wypełnienie $B_p$;
\href{https://arxiv.org/abs/alg-geom/9505018}
{R.\ Fintushel i R.\ J.\ Stern, \emph{Rational blowdowns
of smooth 4-manifolds}} --- konstrukcja i niezmienniki
kontrakcji wymiernej;
\href{https://arxiv.org/abs/dg-ga/9612014}
{R.\ Fintushel i R.\ J.\ Stern, \emph{Knots, Links,
and 4-Manifolds}} --- formuła chirurgii wzdłuż węzła
i przykłady na K3.\par}
